% ===================================================================
%  Minimalbeispiel: Liedblatt – „Hänsel und Gretel“ (Volkslied, gemeinfrei)
%  G-Dur, 2/4-Takt. Erste Strophe unter den Noten (musixlyr),
%  Strophen 2–3 zweispaltig darunter, alles auf einer Seite.
%  Eigenständig: braucht keine header.tex.
%
%  Drei Durchläufe:  pdflatex … → musixflx … → pdflatex …
%  (musixflx gleicht alle Notenzeilen auf dieselbe Länge aus.)
% ===================================================================
\documentclass[a4paper,12pt]{article}
\usepackage[T1]{fontenc}
\usepackage[utf8]{inputenc}
\usepackage[ngerman]{babel}
\usepackage{lmodern}
\usepackage[margin=1.8cm]{geometry}
\usepackage{multicol}       % Strophen zweispaltig
\usepackage{musixtex}
\input musixlyr             % Liedtext unter den Noten
\pagestyle{empty}
\setlength{\columnsep}{1em}

% Eine Strophe: \strophe{Nummer}{Text}; bleibt beim Spaltenumbruch zusammen.
\newcommand{\strophe}[2]{%
  \noindent\makebox[1.4em][l]{\textbf{#1.}}%
  \begin{minipage}[t]{\dimexpr\linewidth-1.4em}\raggedright #2\end{minipage}%
  \par\vspace{1.2em}}

\begin{document}
\begin{center}
{\LARGE\bfseries\sffamily Hänsel und Gretel}\\[4pt]
{\small Text und Melodie: Volksweise \textbullet\ G-Dur}
\end{center}
\vspace{1em}

\Largemusicsize             % großes System: VOR \begin{music}
\begin{music}
\parindent0pt                % erste Zeile nicht einrücken: alle Zeilen gleich lang
\staffbotmarg=8\Interligne  % Platz unter jedem System (innerhalb von music)
\instrumentnumber{1}\setclef1\treble
\generalmeter{\meterfrac24}\generalsignature{1}\nobarnumbers
\setsongraise1{11mm}         % Text dicht unter die eigene Zeile
\setlyrics{strophe}{Hän-sel und Gre-tel ver-lie-fen sich im Wald.
Es war so fins-ter und auch so bit-ter kalt.
Sie ka-men an ein Häus-chen von Pfef-fer-ku-chen fein.
Wer mag der Herr wohl von die-sem Häus-chen sein?}
\assignlyrics1{strophe}
\startpiece
% Tonhöhen: g = g', h = a', i = h', j = c'', k = d''.
% Achtelpaare mit Balken: \ibu (Balken oben) bzw. \ibl (unten), \qb, \tbu/\tbl.
\NOtes\ql k\en\Notes\ibl0i1\qb0i\tbl0\qb0j\en\bar                    % Hän-sel und
\NOtes\ql k\en\Notes\ibu0i{-2}\qb0i\tbu0\qb0g\en\bar                 % Gre-tel ver-
\Notes\ibu0h0\qb0h\tbu0\qb0h\ibu0h1\qb0h\tbu0\qb0i\en\bar            % lie-fen sich im
\NOtes\qu g\qp\en                                                    % Wald.
\alaligne
\NOtes\ql k\en\Notes\ibl0i1\qb0i\tbl0\qb0j\en\bar                    % Es war so
\NOtes\ql k\en\Notes\ibu0i{-2}\qb0i\tbu0\qb0g\en\bar                 % fins-ter und
\Notes\ibu0h0\qb0h\tbu0\qb0h\ibu0h1\qb0h\tbu0\qb0i\en\bar            % auch so bit-ter
\NOtes\qu g\en\Notes\ds\cu g\en                                      % kalt. Sie
\alaligne
\Notes\ibu0h0\qb0h\tbu0\qb0h\ibu0h1\qb0h\tbu0\qb0i\en\bar            % ka-men an ein
\NOtes\ql j\en\Notes\ibu0h0\qb0h\tbu0\qb0h\en\bar                    % Häus-chen von
\Notes\ibu0i{-1}\qb0i\tbu0\qb0h\ibl0i1\qb0i\tbl0\qb0j\en\bar         % Pfef-fer-ku-chen
\NOtes\ql k\qp\en                                                    % fein.
\alaligne
\NOtes\ql k\en\Notes\ibl0i1\qb0i\tbl0\qb0j\en\bar                    % Wer mag der
\NOtes\ql k\en\Notes\ibu0i{-2}\qb0i\tbu0\qb0g\en\bar                 % Herr wohl von
\Notes\ibu0h0\qb0h\tbu0\qb0h\ibu0h1\qb0h\tbu0\qb0i\en\bar            % die-sem Häus-chen
\NOTes\hu g\en                                                       % sein?
\Endpiece
\end{music}
\assignlyrics1{}            % Zuweisung beenden (gilt sonst global weiter)

\vspace{2em}

\begin{multicols}{2}
\strophe{2}{Hu, hu, da schaut\\ eine alte Hexe raus!\\
Sie lockt die Kinder\\ ins Pfefferkuchenhaus.\\
Sie stellte sich gar freundlich,\\ o Hänsel, welche Not!\\
Ihn wollt sie braten\\ im Ofen braun wie Brot.}

\columnbreak
\strophe{3}{Doch als die Hexe\\ zum Ofen schaut hinein,\\
ward sie gestoßen\\ von unserm Gretelein.\\
Die Hexe musste braten,\\ die Kinder gehn nach Haus.\\
Nun ist das Märchen\\ von Hans und Gretel aus.}
\end{multicols}

\end{document}
